\subsection*{Partie A}

On considère la fonction $f$ définie sur l'intervalle $] - 1~;~+\infty[$ par 

\[f(x) = 4\ln(x + 1) - \dfrac{x^2}{25}\]

On admet que la fonction $f$ est dérivable sur l'intervalle $] - 1~;~+\infty[$.

\medskip

\begin{enumerate}
\item Déterminer la limite de la fonction $f$ en $-1$.
\item Montrer que, pour tout $x$ appartenant à l'intervalle $]- 1~;~+\infty[$ , on a :
 \[f'(x) = \dfrac{100- 2x - 2x^2}{25(x+1)}\]

\item Étudier les variations de la fonction $f$ sur l'intervalle $] - 1~;~+\infty[$ puis en déduire que la fonction $f$ est strictement croissante sur l'intervalle [2~;~6,5].
\item On considère $h$ la fonction définie sur l'intervalle [2~;~6,5] par $h(x) = f(x) - x$.

On donne ci-dessous le tableau de variations de la fonction $h$ :

\begin{center}
\psset{unit=1cm,arrowsize=2pt 3}
\begin{pspicture}(11,2.5)
\psframe(11,2.5)\psline(0,2)(11,2) \psline(1,0)(1,2.5)
\uput[u](0.5,1.9){$x$}\uput[u](1.2,1.9){$2$}\uput[u](6,1.9){$m\approx 2,364$}\uput[u](10.75,1.9){$6,5$}
\rput(0.5,1){$h(x)$}\uput[u](1.4,0){$h(2)$}\uput[d](6,2){$M \approx 2,265$}\uput[u](10.5,0){$h(6,5)$}
\psline{->}(1.8,0.4)(4.9,1.6)\psline{->}(6.9,1.6)(10,0.4)
\end{pspicture}
\end{center}
Montrer que l’équation $h(x) = 0$ admet une unique solution $\alpha \in [2~;~6,5]$.

\item On considère le script suivant, écrit en langage Python:

\begin{center}
\begin{tabular}{|l|}\hline
\quad \textbf{from} math \textbf{import *}\\
~\\
\quad \textbf{def} f(x):\\
\quad \qquad\textbf{return} 4*log(1+x)-(x**2)/25\\
~\\
\quad \textbf{def} bornes(n):\\
\quad \qquad p = 1/10**n\\
\quad \qquad x = 6\\
\quad \qquad \textbf{while} f(x)-x > 0:\\
\quad \qquad\qquad x = x + p\\
\quad \qquad\textbf{return} (x-p, x)\\\hline
\end{tabular}
\end{center}

On rappelle qu'en langage Python:


\begin{itemize}[label=$\bullet~$]
\item la commande $\text{log} (x)$ renvoie la valeur $\ln x$ ; 
\item la commande c$**$d renvoie la valeur de $\text{c}^{\text{d}}$.
\end{itemize}

	\begin{enumerate}
		\item Donner les valeurs renvoyées par la commande bornes(2).
		
On donnera les valeurs arrondies au centième.
		\item Interpréter ces valeurs dans le contexte de l'exercice.
	\end{enumerate}
\end{enumerate}

\subsection*{Partie B}

\emph{Dans cette partie, on pourra utiliser les résultats obtenus dans la partie A.}

On considère la suite $(u_n)$ définie par $u_0 = 2$, et, pour tout entier naturel $n,\: u_{n+1} = f\left(u_n\right)$

\begin{enumerate}
\item Montrer par récurrence que pour tout $n$ entier naturel,

\[2 \leqslant u_n \leqslant u_{n+1} < 6,5.\]
\item En déduire que la suite $(u_n)$ converge vers une limite $\ell$.
\item On rappelle que le réel $\alpha$, défini dans la partie A, est la solution de l'équation $h(x)~=~0$ sur l'intervalle [2~;~6,5].

Justifier que $\ell = \alpha$.
\end{enumerate}

\bigskip
